%!TEX TS-program = pdflatex
%!TEX encoding = UTF-8 Unicode
%!TEX spellcheck = English (Aspell)

\documentclass[11pt]{article}
\usepackage{lucy,handout,graphicx}

% =========================================================
\begin{document}

Below are a few graph problems that can be solved by reducing to
computing optimal spanning trees or maximum bottleneck paths. (Recall that
both of these can be done with variants of ``best'' first search!)

Whenever you write an algorithm that reduces to a graph problem, make sure that someone
can answer these questions when they read your solution:

\begin{itemize}
\item What are the vertices?  What does each vertex represent?
\item What are the edges?  Are they directed or undirected?
\item If the vertices and/or edges have associated values, what are they?
\item \textbf{What is the precise relationship between your graph and the stated problem?}
\item What standard graph algorithm are you using to solve the problem?
\item What is the running time of your entire algorithm, \emph{including} the time to build the graph, \emph{as a function of the original input parameters}?
\end{itemize}

And as always, it's great to make up some examples to be sure you understand the problem
you are trying to solve.

\begin{enumerate}
    \item An international conference has delegates speaking $n$ languages. There are $m$ pairs
    of languages for which you can hire translators, but each one has a cost to
    work the conference. A message may
    pass through several translators.

    \begin{enumerate}
        \item Design and analyze an algorithm to decide whether the available translators
        collectively allow communication between all $n$ languages.
        \item Assuming communication is possible, design and analyze an algorithm to choose
        the cheapest set of translators to allow communication between all $n$ languages.
    \end{enumerate}
    
\item

A museum must move a statue from its storage room $s$ to an exhibition room
$t$. Rooms are connected by doorways, each with a known usable width. A statue
of width $x$ can pass through a doorway exactly when the doorway's width is at
least $x$. Assume height, turning space, and all other constraints are
irrelevant.

Design an algorithm to find the widest statue that can be moved from $s$ to
$t$, and a route for moving it.

\end{enumerate}


%----------------------------------------------------------------------


\end{document}